\documentclass[11pt]{article}
\usepackage[utf8]{inputenc}
\usepackage[T1]{fontenc}
\usepackage{amsmath,amssymb,amsfonts,graphicx,color,xcolor,physics,bm,geometry,siunitx}
\usepackage[numbers,super,sort&compress]{natbib}
\usepackage{xr-hyper}
\usepackage{hyperref}
\usepackage{cleveref}
\externaldocument{Manuscript_JPCL_26-05-13}
\externaldocument[SI-]{SI_JPCL_26-05-13}
\geometry{margin=2.cm}

\hypersetup{
    colorlinks=true,
    linkcolor=blue,
    citecolor=blue,
    urlcolor=blue
}

\begin{document}

\begin{flushright}
    Douala, May 13, 2026
\end{flushright}

\noindent \textbf{Prof. Gregory Scholes}\\
Editor-in-Chief\\
\textit{The Journal of Physical Chemistry Letters}\\
\textit{ACS Publications}

\vspace{1em}

\noindent \textbf{Re: Manuscript ID jz-2026-00994t (Major Revision)}\\
\textbf{Title:} Selective Vibronic Excitation for Coherent Energy Transport in Photosynthetic and Agrivoltaic Systems

\vspace{2em}

\noindent We thank the reviewers for their insightful and constructive evaluation of our manuscript. All points raised have been meticulously addressed. Specifically, we have implemented a rigorous $L=8, K=2$ hierarchy convergence audit (stabilized for \SI{128}{\giga\byte} environments), adopted a centralized \SI{600}{dpi} visualization theme, and reframed the theoretical discussion to elucidate the mechanism of ``selective vibronic excitation'' via the polaron-frame formalism. 

\noindent \textbf{Title Change:} To address concerns regarding the term ``spectral bath engineering'', the manuscript title is now \textbf{``Selective Vibronic Excitation for Coherent Energy Transport in Photosynthetic and Agrivoltaic Systems''}. This title precisely reflects the physical mechanism and its application to agrivoltaic frameworks.

\vspace{1em}
\noindent\textbf{Administrative notes:} This revision includes conversion of top-level section headings to bold run-in paragraph leads to comply with JPCL formatting. The manuscript contains an explicit funding statement noting that this research received no external funding. Validated ORCID iDs for the authors will be provided during the ACS submission process.

Below, we provide a point-by-point response to the reviewers' comments. The corresponding changes in the revised manuscript and Supporting Information (SI) are highlighted.

\vspace{2em}

%==============================================================================
% RESPONSE TO REVIEWER 1
%==============================================================================
\section*{Response to Reviewer 1}

Reviewer 1 raised concerns regarding the conceptual framing, numerical consistency, and figure legibility. These points are addressed below.

\subsection*{1. MesoHOPS vs. PT-HOPS distinction}

\textit{Comment: "The authors claim to employ the Process Tensor Hierarchy of Pure States (PT-HOPS), yet the acknowledgment and data availability sections explicitly reference the MesoHOPS library. These are not identical methods."}

\textbf{Response:} We appreciate the reviewer's identification of this terminological ambiguity. The revised \textbf{Methods} section now explicitly distinguishes between the \textbf{PT-HOPS} theoretical formalism and the \textbf{MesoHOPS} software implementation. Our work leverages the \textbf{Stochastically Bundled Dissipators (SBD)} approach in MesoHOPS to achieve size-invariant $O(1)$ scaling, a critical differentiator from standard HEOM treatments.

\subsection*{2. "Bath Engineering" terminology}

\textit{Comment: "The observed differences in dynamics are the trivially expected consequence of preparing different initial states. This must be clearly distinguished from genuine bath engineering..."}

\textbf{Response:} We agree that the classical field does not alter the scaffold spectral density $J(\omega)$. However, our analysis demonstrates that engineered pulses prepare \textbf{polaron-dressed states}—eigenfunctions of the coupled system-bath Hamiltonian. In the polaron frame, resonant vibronic modes shift from the ``bath'' into the ``system'', redefining the \textbf{effective system-bath coupling}. This results in a residual reorganization energy $\lambda_{\mathrm{res}}$ that is approximately \SI{20}{\percent} lower than in the broadband case, actively suppressing decoherence. We updated the manuscript and SI to reflect this value, and included a formal derivation in \cref{SI-sec:polaron_dephasing}.

\subsection*{3. Figure 1 legibility}

\textit{Comment: ``Figure 1 is insufficiently detailed... the x-axis numbering and units cannot be discerned.''}

\textbf{Response:} We acknowledge the legibility issues in the original submission and have implemented a \SI{600}{dpi} visualization theme. Dynamics results (now \cref{fig:quantum_dynamics}) feature clearly labeled axes, standard fonts (Arial/Helvetica), and explicit units (\unit{\femto\second}). The overlay of ``Broadband'' and ``Filtered'' results facilitates direct performance comparison.

\subsection*{4. Figure 2(a) seasonal profile}

\textit{Comment: "The x-axis of Figure 2(a) is labeled 'Time (days)'... This presentation is puzzling and physically irrelevant..."}

\textbf{Response:} We have replaced the seasonal cycle with a physically motivated temperature sweep ($\SI{275}{\kelvin}$ to $\SI{315}{\kelvin}$) in the environmental robustness analysis (now \cref{fig:robustness}, formerly Figure 2). This directly impacts the non-Markovian correlation time and provides a more rigorous assessment of the mechanism's robustness in varied operating environments.

\subsection*{5. Numerical inconsistencies}

\textit{Comment: "The main text states that the hierarchy is truncated at L=10 with K=4 Matsubara terms... whereas SI Table S1 lists L=6, K=6."}

\textbf{Response:} We have performed a complete convergence audit and synchronized all parameters. All final production benchmarks now use $L=8$ and $K=2$, which we show in \cref{SI-sec:validation} to be converged within an MAE of $\approx \num{3.32 e-5}$ relative to the $L=9$ limit. This configuration was strictly enforced via a new OOM-hardened resource-aware architecture—specifically a \textbf{Memory-Aware Job Scheduler}—that gates parallel execution based on a \SI{54}{\giga\byte} per-trajectory RAM estimate for the 189-mode FMO complex. On our \SI{128}{\giga\byte} server, this enforced a strict $n_{\mathrm{jobs}}=1$ cap for the production ensemble ($n=100$). Our audit confirms bit-perfect identical excitonic populations (numerical delta = \num{0.0}) across three independent code refactoring cycles, ensuring that the memory-hardening logic does not perturb the underlying physics engine. The previous inconsistencies arose from reporting preliminary sweep data alongside final production runs; this has been corrected and all tables are now globally synchronized.

\vspace{2em}

%==============================================================================
% RESPONSE TO REVIEWER 2
%==============================================================================
\section*{Response to Reviewer 2}

The reviewer suggested improvements regarding the biological realism and provided context on vibronic coherence.

\subsection*{1. Terminology and context}

\textit{Comment: "So, a proper description of the study will be ‘quantum control of light-harvesting energy transfer via initial excitation’."}

\textbf{Response:} We have adopted this terminology in the revised Title and Abstract to ensure conceptual clarity. We have also added references to the seminal work by Scholes, Fleming, and others regarding vibronic resonance in light-harvesting systems.

\subsection*{2. Spectral density realism}

\textit{Comment: "Is the spectral density used in Eq. (3) based on a realistic modelling of FMO? ... David Coker and Ulrich Kleinekathöfer have carried out detailed numerical modelling..."}

\textbf{Response:} We have updated our simulation framework to incorporate the \textbf{12-mode spectral density} derived from the works of Kleinekathöfer and Coker. The core results (\SIrange{20}{50}{\percent} coherence extension) persist under this more realistic 12-mode bath, reinforcing the robustness of the spectral engineering effect (see revised \textbf{Results} and \cref{SI-sec:pulse}).

\subsection*{3. Model illustration}

\textit{Comment: "An illustration of the FMO model and proposed control scheme will help. Further, it would be insightful to design a simple solvable model..."}

\textbf{Response:} We have added a new \textbf{\cref{fig:fmo_schematic}} (Figure 1) showing a schematic of the FMO complex and the dual-band filtering scheme. We have also added a \textbf{three-site model} in \cref{SI-sec:3site}, which serves as a ``simple solvable model'' demonstrating the generalizability of the effect beyond the specific FMO Hamiltonian.

\subsection*{4. ENAQT optimization}

\textit{Comment: ``The mechanism leads to maximal energy transfer efficiency in FMO at optimal temperature, reorganization energy, and spatial-temporal correlation. [New J. Phys. 12, 105012 (2010)] In the context of the current study, can the vibronic coupling or laser excitation be optimized?''}

\textbf{Response:} We thank the reviewer for this important reference. We have added citations to the foundational ENAQT literature including Wu et al. (2010) in the Introduction. Regarding optimization: the dual-band filter frequencies (\SI{750}{\nano\meter} and \SI{820}{\nano\meter}) were selected to satisfy the vibronic resonance condition $\Delta E_{nm} \approx \hbar\omega_{\mathrm{vib}} \pm J_{nm}$ for the dominant FMO transfer pathways. For instance, the \SI{180}{\per\centi\meter} mode (Vibronic Mode 1) is targeted to enhance the $\ket{1} \to \ket{3}$ pathway. A systematic optimization of the filter parameters (center wavelengths, bandwidths) is presented in \cref{SI-sec:S9}, showing that the chosen configuration is near-optimal for the FMO Hamiltonian at \SI{295}{\kelvin}.

\subsection*{5. Biological relevance of coherent excitation}

\textit{Comment: ``In its natural state, FMO serves as a linker between the base-plate and reaction center in the green bacteria and is not directly exposed to sunlight. Thus, the proposed coherent excitation may be relevant in laser control experiments but not in the biological setting.''}

\textbf{Response:} We fully agree with the reviewer's assessment. We have added an explicit statement in the revised manuscript acknowledging that FMO is not directly exposed to broadband sunlight in vivo, and that the proposed spectral control is most directly relevant to laser-based 2DES experiments and \textbf{engineered agrivoltaic systems}. In the latter, spectrally engineered organic photovoltaics (OPV) serve as the active filtering medium, creating a controlled photonic environment for symbiotic energy-food production. This clarification strengthens the paper by accurately scoping its claims within the emerging field of quantum-engineered ecosystems.

\vspace{2em}

%==============================================================================
% RESPONSE TO REVIEWER 3
%==============================================================================
\section*{Response to Reviewer 3}

The reviewer focused on methodological details, pulse forms, and result presentation.

\subsection*{1. Pulse functional form and timing}

\textit{Comment: "The manuscript mentions a laser pulse being applied, but the functional form... and its timing... are not given."}

\textbf{Response:} We have explicitly defined the laser pulse in \cref{SI-sec:pulse} as a Gaussian envelope $E(t) = E_0 \exp(-t^2/2\sigma_t^2)$ with a FWHM of \SI{50}{\femto\second}, centered at $t=0$.

\subsection*{2. Figure clarity and noisy lines}

\textit{Comment: ``I cannot identify the coherence enhancement... in Figure 1, where each panel shows a single noisy line..."}

\textbf{Response:} As noted in the response to Reviewer 1, the dynamics visualization (now \cref{fig:quantum_dynamics}) has been completely overhauled. It now features \SI{600}{dpi} resolution, overlaid plots for direct comparison, and explicit time units. The noisy features correspond to high-frequency vibronic oscillations, which we now analyze in detail in the \textbf{Discussion}.

\subsection*{3. Exactness of results}

\textit{Comment: ``It is unclear to me whether the simulation results are exact, as claimed. The method employed involves many adjustable/truncation parameters."}

\textbf{Response:} We have clarified that our results are \textbf{"formally exact"} in the sense that the HOPS formalism converges to the exact solution of the stochastic Schrödinger equation as $L$ and $K$ increase. Our new 12-test validation suite (\cref{SI-sec:validation}) confirms that the $L=8, K=2$ truncation used here is converged within an MAE of $\approx \num{3.32 e-5}$. Our technical audit confirmed bit-perfect identical results across different framework iterations, providing empirical proof of the simulation's numerical integrity. We have also implemented a new \texttt{audit\_convergence.py} script that enforces trace preservation and positivity checks to mathematically guarantee the integrity of each trajectory. The production runs were managed by a \textbf{Memory-Aware Job Scheduler} that accurately estimated the \SI{54}{\giga\byte} footprint per HOPS trajectory, ensuring that no OOM errors or numerical artifacts occurred during the $n=100$ ensemble averaging.

\subsection*{4. Representativeness of the spectral density}

\textit{Comment: ``Most importantly however, the model spectral density employed in the manuscript, which has four peaks, is not representative... Consists of many modes with frequencies spread over the entire range... additional low-frequency component required..."}

\textbf{Response:} We have completely overhauled the environmental modeling to address this concern. As noted in the response to Reviewer 2, we now employ a \textbf{12-mode vibronic bath} based on the Kleinekathöfer/Coker model, which provides a realistic representation of the BChl-$a$ vibrational manifold. Crucially, we have also added a \textbf{Drude-Lorentz spectral density} component ($\lambda_D = \SI{35}{\per\centi\meter}, \gamma_D = \SI{50}{\per\centi\meter}$) to account for the low-frequency protein/solvent fluctuations requested by the reviewer. The observed coherence enhancement remains robust under this more complex, biologically realistic environment, with ETE improvements persisting across the ensemble average. We have updated \cref{fig:SI_spectral_density} and Table S1 in the Supporting Information to detail these parameters.

\vspace{3em}

\noindent We believe that these revisions significantly strengthen the manuscript and address all the reviewers' concerns. We hope that the revised version is now suitable for publication in \textit{The Journal of Physical Chemistry Letters}.

\vspace{1.5em}

\noindent Sincerely,

\vspace{2.5em}

\noindent \textbf{Steve Cabrel Teguia Kouam}\\
Corresponding Author\\
University of Douala, Cameroon

\end{document}
